\begin{tikzpicture}[
    section/.style={
        draw,
        semithick,
        align=center,
        minimum height=14mm,
        minimum width=28mm,
        text width=24mm,
        outer sep=0pt,
        font=\small
    },
    header/.style={section, fill=gray!12},
    block/.style={section, fill=gray!12},
    codeblock/.style={section, fill=gray!12},
    field/.style={
        draw,
        semithick,
        align=center,
        minimum height=12mm,
        minimum width=32mm,
        text width=29mm,
        outer sep=0pt,
        fill=gray!12,
        font=\scriptsize
    },
    detail/.style={
        draw,
        semithick,
        align=center,
        minimum height=9mm,
        minimum width=34mm,
        text width=31mm,
        outer sep=0pt,
        fill=gray!12,
        font=\scriptsize
    },
    widefield/.style={
        draw,
        semithick,
        align=center,
        minimum height=12mm,
        minimum width=45mm,
        text width=42mm,
        outer sep=0pt,
        fill=gray!12,
        font=\scriptsize
    },
    codefield/.style={widefield, fill=green!17},
    note/.style={
        draw,
        semithick,
        align=center,
        rounded corners=2pt,
        minimum height=9mm,
        minimum width=64mm,
        text width=60mm,
        fill=gray!12,
        font=\scriptsize
    },
    attach/.style={->, semithick},
    weakattach/.style={->, semithick}
]

\path coordinate (mainaux) at (0,0)
    node[
        header,
        anchor=west,
        label={[name=main-l1, align=center, font=\small, text width=31mm]center:
            {заголовок\\24 байта}}
    ] (header) at (mainaux) {}
    node[
        block,
        anchor=west,
        label={[name=main-l2, align=center, font=\footnotesize, text width=31mm]center:
            {таблица строк\\переменный размер}}
    ] (strings) at (header.east) {}
    node[
        section,
        fill=green!17,
        anchor=west,
        label={[name=main-l3, align=center, font=\small, text width=31mm]center:
            {импорты\\$n_i \times 4$ байта}}
    ] (imports) at (strings.east) {}
    node[
        section,
        fill=green!17,
        anchor=west,
        label={[name=main-l4, align=center, font=\footnotesize, text width=31mm]center:
            {публичные \\ символы\\$n_p \times 9$ байт}}
    ] (publics) at (imports.east) {}
    node[
        codeblock,
        anchor=west,
        minimum width = 12mm,
        label={[name=main-l5, align=center, font=\small, text width=16mm]center:
            {байткод \\ до \texttt{0xFF}}}
    ] (code) at (publics.east) {};

\node[detail, fill=green!17] (h0) at ($(header.west)+(-2.8,2.25)$)
{магическое число: \\ \texttt{uint8[4]}};
\node[detail, fill=green!17, below=0pt of h0] (h1)
{версия: \texttt{int32}};
\node[detail, below=0pt of h1] (h2)
{размер таблицы \\ строк: \texttt{int32}};
\node[detail, below=0pt of h2] (h3)
{число глобальных \\ переменных: \texttt{int32}};
\node[detail, fill=green!17, below=0pt of h3] (h4)
{число импортов: \texttt{int32}};
\node[detail, below=0pt of h4] (h5)
{число публичных\\символов: \texttt{int32}};
\draw[attach] (header.west) -- (h3.east);

\path coordinate (impaux) at ($(imports.south)+(-3.2,-1.65)$)
    node[
        field,
        fill=green!17,
        anchor=west,
        label={[name=imp-label, align=center, font=\scriptsize, text width=42mm]center:
            {смещение имени юнита\\в таблице строк\\\texttt{int32}}}
    ] (impentry) at (impaux) {};
\draw[attach] (imports.south) to[out=-90,in=90] (impentry.north);

\path coordinate (pubaux) at ($(publics.north)+(-6.8,1.8)$)
    node[
        field,
        anchor=west,
        label={[name=pub-l1, align=center, font=\scriptsize, text width=29mm]center:
            {смещение имени\\ в таблице строк \\\texttt{int32}}}
    ] (pub1) at (pubaux) {}
    node[
        field,
        anchor=west,
        label={[name=pub-l2, align=center, font=\scriptsize, text width=29mm]center:
            {адрес функции\\или индекс\\\texttt{int32}}}
    ] (pub2) at (pub1.east) {}
    node[
        field,
        fill=green!17,
        anchor=west,
        label={[name=pub-l3, align=center, font=\scriptsize, text width=29mm]center:
            {флаг\\\texttt{uint8}}}
    ] (pub3) at (pub2.east) {};
\draw[attach] (publics.north) to[out=90,in=-90] (pub2.south);

\end{tikzpicture}
